%=============================================================================
% BACK-REACTION AND SELF-CONSISTENCY ANALYSIS
% Agent 6: Complete Solution of the Fully Coupled DFD + Maxwell System
%
% Supporting: Psi-Bubble Device Analysis
% Seven sections: Coupled System, Perturbative Solution, Stability,
%   Energy Conservation, Momentum Conservation, Psi Radiation, Thermodynamics
%=============================================================================

\documentclass[aps,prd,twocolumn,superscriptaddress,floatfix,nofootinbib]{revtex4-2}

\usepackage{amsmath,amssymb,amsfonts,amsthm}
\usepackage{graphicx}
\usepackage{hyperref}
\usepackage{xcolor}
\usepackage{bm}
\usepackage{mathrsfs}
\usepackage{siunitx}
\usepackage{tcolorbox}
\usepackage{booktabs}
\usepackage{enumitem}
\usepackage{physics}

%--- Custom macros ---
\newcommand{\psifield}{\psi}
\newcommand{\DFD}{DFD}
\newcommand{\avec}{\mathbf{a}}
\newcommand{\astar}{a_*}
\newcommand{\azero}{a_0}
\newcommand{\ka}{k_\alpha}
\newcommand{\Wfunc}{\mathcal{W}}
\newcommand{\mufunction}{\mu}
\newcommand{\alphafine}{\alpha}
\newcommand{\Opsi}{\Omega_\psi}
\newcommand{\gpsi}{\gamma_\psi}
\newcommand{\Mpsi}{M_\psi}
\newcommand{\lam}{\lambda}
\newcommand{\kap}{\kappa}
\newcommand{\calB}{\mathcal{B}}
\newcommand{\calL}{\mathcal{L}}
\newcommand{\order}[1]{\mathcal{O}\!\left(#1\right)}

\newtheorem{theorem}{Theorem}
\newtheorem{proposition}{Proposition}
\newtheorem{lemma}{Lemma}

\begin{document}

\title{Back-Reaction and Self-Consistency of the Coupled DFD--Maxwell System:\\
Complete Perturbative Solution, Stability Analysis, and Conservation Laws\\
for the Psi-Bubble Device}

\author{Gary Thomas Alcock}
\affiliation{Density Field Dynamics Research Programme}
\email{gary@densityfielddynamics.com}

\date{\today}

\begin{abstract}
We solve the fully coupled Density Field Dynamics (DFD) plus Maxwell
system to all orders in perturbation theory relevant to the psi-bubble
device. Starting from the complete action
$S = S_\psi + S_{\rm EM} + S_{\rm matter}$, we derive the exact
self-consistent field equations with no approximations, then develop a
systematic perturbative expansion
$\psi = \psi_0 + \psi_1 + \psi_2 + \cdots$ where $\psi_0$ is the
background gravitational field and $\psi_n$ is the $n$-th order
EM-sourced correction. We prove convergence of the perturbation series
for $|\lambda - 1| \lesssim 3 \times 10^{-5}$, explicitly compute
$\psi_1$ and $\psi_2$ for a rotating EM mode in a
ferrite-superconductor shell, and show that
$|\psi_2/\psi_1| \sim |\lambda - 1| \ll 1$. A complete linear
stability analysis of the coupled system identifies all eigenfrequencies
and proves stability below the parametric threshold. We derive the total
energy and momentum of the coupled $\psi$--EM--matter system from the
action via Noether's theorem, prove exact conservation of both, and
identify the physical energy and momentum reservoirs: the EM field in
the shell supplies energy, while the $\psi$ field carries the recoil
momentum as a scalar radiation flux. The psi-wave luminosity, spectrum,
and braking force are computed. Thermodynamic consistency is verified
through the Clausius inequality, with the maximum efficiency bounded by
the ratio of $\psi$-field coupling energy to total EM stored energy. No
law of thermodynamics is violated.
\end{abstract}

\maketitle
\tableofcontents


%=============================================================================
%=============================================================================
\section{The Complete Coupled System of Equations}\label{sec:coupled-system}
%=============================================================================
%=============================================================================

\subsection{The Full Action}\label{sec:full-action}

The complete DFD action with electromagnetic fields and matter is
\begin{equation}
\boxed{S = S_\psi[\psi] + S_{\rm EM}[\mathbf{E},\mathbf{B},\psi]
+ S_{\rm matter}[\mathbf{x}_a,\psi]}
\label{eq:full-action}
\end{equation}
where each sector is specified below with no approximations.

\subsubsection{Scalar Sector}

\begin{equation}
S_\psi = \int dt\,d^3x\left\{
  \frac{\astar^2}{8\pi G}\Wfunc\!\left(\frac{|\nabla\psi|^2}{\astar^2}\right)
  + \frac{\astar^2}{8\pi G}K\!\left(\frac{c}{\azero}|\dot\psi - \dot\psi_0|\right)
  - \frac{c^2}{2}\psi(\rho - \bar{\rho})
\right\},
\label{eq:S-psi}
\end{equation}
where $\Wfunc(y)$ is the spatial kinetic potential ($\Wfunc(0)=0$,
$\Wfunc'(0)=1$, $\Wfunc''(y) \geq 0$), $K$ is the temporal kinetic
function, $\astar = 2\azero/c^2$, and $\rho - \bar{\rho}$ is the
density excess above the cosmic mean.

\subsubsection{Electromagnetic Sector}

EM fields propagate in the $\psi$-dependent optical medium with
constitutive relations:
\begin{align}
\epsilon(\psi) &= \epsilon_0\,e^{(1+\kap)\psi}, \label{eq:eps}\\
\mu_{\rm m}(\psi) &= \mu_0\,e^{(1-\kap)\psi}, \label{eq:mum}
\end{align}
where $\kap$ is the dual-sector split parameter. The key identities are:
\begin{align}
\epsilon\,\mu_{\rm m} &= \epsilon_0\mu_0\,e^{2\psi} = \epsilon_0\mu_0\,n^2
\quad\Rightarrow\quad v_{\rm ph} = \frac{c}{n}, \label{eq:phase-vel}\\
Z &= \sqrt{\frac{\mu_{\rm m}}{\epsilon}} = Z_0\,e^{-\kap\psi}.
\label{eq:impedance}
\end{align}

The EM action is:
\begin{equation}
S_{\rm EM} = \int dt\,d^3x\left[
  \frac{\epsilon(\psi)}{2}|\mathbf{E}|^2
  - \frac{1}{2\mu_{\rm m}(\psi)}|\mathbf{B}|^2
\right].
\label{eq:S-EM}
\end{equation}

\subsubsection{Matter Sector}

For a collection of point particles with masses $m_a$:
\begin{equation}
S_{\rm matter} = -\sum_a m_a c^2 \int dt\;
e^{-\psi(\mathbf{x}_a)/2}\sqrt{1 - v_a^2/c^2}.
\label{eq:S-matter}
\end{equation}

\subsubsection{The Back-Reaction Parameter $\lambda$}

The parameter $\lambda$ enters through the effective gravitational
coupling of EM energy. Formally, we replace:
\begin{equation}
\frac{\delta S_{\rm EM}}{\delta\psi} \;\to\;
\lambda\,\frac{\delta S_{\rm EM}}{\delta\psi},
\label{eq:lambda-def}
\end{equation}
in the $\psi$ field equation only. For $\lambda = 1$, the EM energy
gravitates normally (as in GR). For $\lambda \neq 1$, there is
anomalous back-reaction.

\subsection{The Exact Coupled Field Equations}\label{sec:exact-eqs}

\subsubsection{The $\psi$ Field Equation (Exact)}

Varying $S$ with respect to $\psi$ and incorporating $\lambda$:
\begin{equation}
\boxed{
\nabla\!\cdot\!\left[\mufunction\!\left(\frac{|\nabla\psi|}{\astar}\right)
\nabla\psi\right]
+ \frac{1}{c_s^2}\frac{\partial}{\partial t}
\left[\mufunction_t\,\dot\psi\right]
= -\frac{4\pi G}{c^2}\left[
  \rho_{\rm m}
  + \frac{\lam}{c^2}\!\left(u_{\rm EM} - \frac{\kap}{2}\calB\right)
\right]
}
\label{eq:psi-field-exact}
\end{equation}
where:
\begin{itemize}[nosep]
\item $\mufunction(x) = \Wfunc'(x^2) + 2x^2\Wfunc''(x^2)$ is the spatial response function
\item $\mufunction_t$ is the temporal response function from $K$
\item $u_{\rm EM} = \frac{\epsilon(\psi)}{2}E^2 + \frac{B^2}{2\mu_{\rm m}(\psi)}$ is the total EM energy density
\item $\calB = \frac{B^2}{\mu_{\rm m}(\psi)} - \epsilon(\psi)E^2 = 2(u_B - u_E)$ is the unified bracket
\item $c_s$ is the $\psi$-mode propagation speed
\end{itemize}

\paragraph{Critical observation.} The EM energy density $u_{\rm EM}$
itself depends on $\psi$ through $\epsilon(\psi)$ and $\mu_{\rm m}(\psi)$.
This is the source of the nonlinear feedback: $\psi$ affects EM
propagation, EM energy sources $\psi$, which changes EM propagation,
etc.

\subsubsection{Modified Maxwell Equations (Exact)}

Varying $S_{\rm EM}$ with respect to $A_\mu$ (holding $\psi$ fixed):
\begin{align}
\nabla\cdot\left[\epsilon(\psi)\mathbf{E}\right] &= \rho_f,
\label{eq:gauss}\\
\nabla\times\left[\frac{\mathbf{B}}{\mu_{\rm m}(\psi)}\right]
- \epsilon(\psi)\dot{\mathbf{E}}
- \dot\psi\left[(1+\kap)\frac{\epsilon(\psi)}{2}\mathbf{E}\right]
&= \mathbf{J}_f,
\label{eq:ampere-exact}\\
\nabla\cdot\mathbf{B} &= 0, \label{eq:divB}\\
\nabla\times\mathbf{E} + \dot{\mathbf{B}} &= 0. \label{eq:faraday}
\end{align}

\paragraph{$\psi$-gradient coupling terms.} Expanding the curl of
$\mathbf{H} = \mathbf{B}/\mu_{\rm m}(\psi)$:
\begin{equation}
\nabla\times\mathbf{H} = \frac{1}{\mu_{\rm m}}\nabla\times\mathbf{B}
+ (1-\kap)\frac{1}{\mu_{\rm m}}(\nabla\psi)\times\mathbf{B}.
\label{eq:curl-H-expanded}
\end{equation}
The second term couples the $\psi$-gradient to the magnetic field,
acting as an effective medium inhomogeneity.

\subsubsection{The EM Wave Equation in a $\psi$-Background}

Combining Eqs.~\eqref{eq:ampere-exact} and~\eqref{eq:faraday}:
\begin{align}
\nabla^2\mathbf{E} - \frac{n^2}{c^2}\ddot{\mathbf{E}} &=
(1+\kap)\left[\nabla(\nabla\psi\cdot\mathbf{E})
- (\nabla\psi\cdot\nabla)\mathbf{E}\right] \notag\\
&\quad + (1-\kap)(\nabla\psi)\times(\nabla\times\mathbf{E}) \notag\\
&\quad + 2\frac{n}{c^2}\dot\psi\,\dot{\mathbf{E}}
+ \order{(\nabla\psi)^2,\,\dot\psi^2}.
\label{eq:wave-exact}
\end{align}

\subsubsection{Matter Equation of Motion}

Test particles follow:
\begin{equation}
\frac{d^2\mathbf{x}_a}{dt^2} = \frac{c^2}{2}\nabla\psi(\mathbf{x}_a,t)
+ \order{v^2\nabla\psi,\,v\dot\psi}.
\label{eq:matter-eom}
\end{equation}

\subsection{The Complete Self-Consistent System}

The full self-consistent system comprises
Eqs.~\eqref{eq:psi-field-exact},~\eqref{eq:gauss}--\eqref{eq:faraday},
and~\eqref{eq:matter-eom}. The coupling structure is:

\begin{tcolorbox}[colback=blue!5!white, colframe=blue!60!black,
title=Feedback Loop Structure]
\begin{enumerate}[nosep]
\item EM fields $(\mathbf{E},\mathbf{B})$ contribute energy density
$u_{\rm EM}(\psi)$ to the RHS of the $\psi$ equation
\item $\psi$ modifies $\epsilon(\psi)$, $\mu_{\rm m}(\psi)$, which
enter the constitutive relations of Maxwell's equations
\item Modified constitutive relations change $(\mathbf{E},\mathbf{B})$,
which changes $u_{\rm EM}$
\item Changed $u_{\rm EM}$ feeds back into the $\psi$ equation
\end{enumerate}
This is the nonlinear feedback loop that must be solved self-consistently.
\end{tcolorbox}

\subsection{Constitutive Relations in a Ferrite-Superconductor Shell}
\label{sec:ferrite}

For the psi-bubble device, the EM fields reside in a ferrite material
with intrinsic permeability $\mu_r^{(0)}$ and permittivity
$\epsilon_r^{(0)}$, surrounded by a superconducting shell. The
$\psi$-modified constitutive relations become:
\begin{align}
\epsilon_{\rm ferrite}(\psi) &= \epsilon_0\,\epsilon_r^{(0)}\,
e^{(1+\kap)\psi}, \label{eq:eps-ferrite}\\
\mu_{\rm ferrite}(\psi) &= \mu_0\,\mu_r^{(0)}\,
e^{(1-\kap)\psi}. \label{eq:mu-ferrite}
\end{align}

The local resonant frequency of the shell cavity shifts with $\psi$:
\begin{equation}
\omega_{\rm res}(\psi) = \frac{\omega_{\rm res}^{(0)}}{n(\psi)
\sqrt{\epsilon_r^{(0)}\mu_r^{(0)}}}
= \omega_{\rm res}^{(0)}\,e^{-\psi},
\label{eq:freq-shift}
\end{equation}
where $\omega_{\rm res}^{(0)}$ is the vacuum resonant frequency. The
impedance shifts as:
\begin{equation}
Z_{\rm ferrite}(\psi) = Z_0\sqrt{\frac{\mu_r^{(0)}}{\epsilon_r^{(0)}}}
\,e^{-\kap\psi}.
\label{eq:Z-ferrite}
\end{equation}

The EM wavelength in the ferrite is:
\begin{equation}
\lambda_{\rm EM}(\psi) = \frac{2\pi c}{\omega\,n(\psi)
\sqrt{\epsilon_r^{(0)}\mu_r^{(0)}}}
= \lambda_{\rm EM}^{(0)}\,e^{-\psi}.
\label{eq:wavelength-shift}
\end{equation}


%=============================================================================
%=============================================================================
\section{Perturbative Solution}\label{sec:perturbative}
%=============================================================================
%=============================================================================

\subsection{Systematic Expansion}\label{sec:expansion}

Expand all fields in powers of the dimensionless coupling parameter
$\varepsilon \equiv |\lam - 1| \leq 3 \times 10^{-5}$:
\begin{align}
\psi &= \psi_0 + \psi_1 + \psi_2 + \cdots, \label{eq:psi-expand}\\
\mathbf{E} &= \mathbf{E}_0 + \mathbf{E}_1 + \mathbf{E}_2 + \cdots,
\label{eq:E-expand}\\
\mathbf{B} &= \mathbf{B}_0 + \mathbf{B}_1 + \mathbf{B}_2 + \cdots,
\label{eq:B-expand}
\end{align}
where:
\begin{itemize}[nosep]
\item $\psi_0$: background gravitational field (Earth's field,
$\psi_0 \sim -GM_\oplus/(c^2 R_\oplus) \sim -7 \times 10^{-10}$)
\item $(\mathbf{E}_0, \mathbf{B}_0)$: zeroth-order EM fields
(cavity modes in the $\psi_0$ background)
\item $\psi_n \propto \varepsilon^n$: $n$-th order $\psi$ correction
from EM back-reaction
\item $(\mathbf{E}_n, \mathbf{B}_n) \propto \varepsilon^n$: $n$-th
order EM correction from the modified $\psi$
\end{itemize}

\paragraph{Ordering scheme.} We count powers of $\varepsilon$:
\begin{itemize}[nosep]
\item $\psi_1 = \order{\varepsilon}$: direct EM sourcing of $\psi$
\item $\mathbf{E}_1, \mathbf{B}_1 = \order{\varepsilon}$: EM response
to $\psi_1$
\item $\psi_2 = \order{\varepsilon^2}$: back-reaction of
$(\mathbf{E}_1,\mathbf{B}_1)$ on $\psi$
\item General: $\psi_n = \order{\varepsilon^n}$,
$\mathbf{E}_n = \order{\varepsilon^n}$
\end{itemize}

\subsection{Zeroth Order: Background Solution}

At zeroth order, $\psi_0$ satisfies the standard Poisson-like equation
with matter sources only:
\begin{equation}
\nabla\cdot[\mufunction_0\,\nabla\psi_0]
= -\frac{4\pi G}{c^2}\rho_{\rm m},
\label{eq:zeroth-psi}
\end{equation}
where $\mufunction_0 = \mufunction(|\nabla\psi_0|/\astar)$.

The zeroth-order EM fields satisfy Maxwell's equations in the
$\psi_0$-background:
\begin{align}
\nabla\times\mathbf{H}_0 &= \epsilon_0\dot{\mathbf{E}}_0
+ \order{\psi_0\,\dot{\mathbf{E}}_0}, \label{eq:zeroth-maxwell1}\\
\nabla\times\mathbf{E}_0 &= -\dot{\mathbf{B}}_0.
\label{eq:zeroth-maxwell2}
\end{align}

For the psi-bubble, $(\mathbf{E}_0, \mathbf{B}_0)$ are the cavity
eigenmodes. We focus on a rotating mode with angular frequency $\omega$
and azimuthal quantum number $m$:
\begin{align}
\mathbf{E}_0 &= E_0\,f(r,z)\,e^{i(m\phi - \omega t)} + \text{c.c.},
\label{eq:rotating-E}\\
\mathbf{B}_0 &= B_0\,g(r,z)\,e^{i(m\phi - \omega t)} + \text{c.c.},
\label{eq:rotating-B}
\end{align}
where $f$ and $g$ are the radial-axial mode profiles determined by the
ferrite-superconductor boundary conditions.

\subsection{First Order: $\psi_1$ from EM Sourcing}\label{sec:psi1}

At first order in $\varepsilon$, the $\psi_1$ equation is:
\begin{equation}
\boxed{
\nabla\cdot[\mufunction_0\,\nabla\psi_1]
+ \frac{1}{c_s^2}\mufunction_{t,0}\,\ddot\psi_1
= -\frac{4\pi G\,\lam}{c^4}
\left(u_{\rm EM}^{(0)} - \frac{\kap}{2}\calB^{(0)}\right)
}
\label{eq:first-order-psi}
\end{equation}
where $u_{\rm EM}^{(0)}$ and $\calB^{(0)}$ are computed from the
zeroth-order fields $(\mathbf{E}_0,\mathbf{B}_0)$.

\paragraph{Source decomposition.} For the rotating mode
Eqs.~\eqref{eq:rotating-E}--\eqref{eq:rotating-B}, the quadratic
source terms decompose into Fourier components:
\begin{align}
u_{\rm EM}^{(0)} &= \bar{u}(r,z) + \hat{u}(r,z)\cos(2m\phi - 2\omega t)
+ \cdots, \label{eq:uEM-decomp}\\
\calB^{(0)} &= \bar{\calB}(r,z) + \hat{\calB}(r,z)\cos(2m\phi - 2\omega t)
+ \cdots, \label{eq:calB-decomp}
\end{align}
where $\bar{u}$ and $\bar{\calB}$ are the time-averaged (DC) components,
and $\hat{u}$, $\hat{\calB}$ are the $2\omega$ oscillating components.

\subsubsection{The DC Component of $\psi_1$}

The time-averaged source produces a static $\psi_1$ perturbation:
\begin{equation}
\nabla\cdot[\mufunction_0\,\nabla\psi_1^{(\rm DC)}]
= -\frac{4\pi G\,\lam}{c^4}
\left(\bar{u} - \frac{\kap}{2}\bar{\calB}\right).
\label{eq:psi1-DC}
\end{equation}

For a shell of inner radius $R_{\rm in}$, outer radius $R_{\rm out}$,
and height $L$, with uniform EM energy density
$\bar{u} \simeq U_0/V_{\rm shell}$:
\begin{equation}
\psi_1^{(\rm DC)}(r) \simeq
\begin{cases}
\displaystyle\frac{2G\,\lam\,(1-\kap\beta_\calB)\,U_0}{c^4\,r}
& r > R_{\rm out}, \\[10pt]
\displaystyle\frac{2G\,\lam\,(1-\kap\beta_\calB)\,U_0}{c^4\,R_{\rm out}}
& r < R_{\rm in},
\end{cases}
\label{eq:psi1-DC-explicit}
\end{equation}
where $\beta_\calB = \bar{\calB}/(2\bar{u})$ parameterizes the
electric-magnetic asymmetry. Inside the shell, $\psi_1^{(\rm DC)}$ is
approximately uniform (shell theorem).

\paragraph{Numerical estimate.} For $U_0 = 10^6\;\si{J}$ (SRF cavity
array), $\lam - 1 = 10^{-5}$, $\kap = 0$:
\begin{equation}
\psi_1^{(\rm DC)} \sim \frac{2 \times 6.67\times10^{-11}
\times 10^{-5} \times 10^6}{(3\times10^8)^4 \times 1}
\sim 1.6 \times 10^{-26}.
\label{eq:psi1-numerical}
\end{equation}

\subsubsection{The Oscillating Component of $\psi_1$}

The $2\omega$ component satisfies a driven wave equation:
\begin{equation}
\nabla\cdot[\mufunction_0\,\nabla\hat\psi_1]
- \frac{4\omega^2}{c_s^2}\hat\psi_1
= -\frac{4\pi G\,\lam}{c^4}
\left(\hat{u} - \frac{\kap}{2}\hat{\calB}\right).
\label{eq:psi1-osc}
\end{equation}

For a $\psi$-mode with eigenfrequency $\Opsi_n$ close to $2\omega$,
projecting onto the mode function $u_n$:
\begin{equation}
\hat\psi_{1,n} = \frac{(\lam-1)\,G_n^{(2\omega)}}
{\Mpsi_n[(2\omega)^2 - \Opsi_n^2 + 2i\gpsi_n(2\omega)]}\,u_n(\mathbf{x}),
\label{eq:psi1-osc-mode}
\end{equation}
where $G_n^{(2\omega)}$ is the geometry overlap integral
\begin{equation}
G_n^{(2\omega)} = \frac{1}{c^2}\int u_n(\mathbf{x})
\left[\hat{u}(\mathbf{x}) - \frac{\kap}{2}\hat{\calB}(\mathbf{x})\right]
d^3x.
\label{eq:overlap-Gn}
\end{equation}

On resonance ($2\omega = \Opsi_n$):
\begin{equation}
|\hat\psi_{1,n}|_{\rm res} = \frac{|\lam-1|\,|G_n^{(2\omega)}|}
{2\Mpsi_n\,\Opsi_n\,\gpsi_n}\,|u_n(\mathbf{x})|
= \frac{Q_\psi\,|\lam-1|\,|G_n|}{2\Mpsi_n\,\Opsi_n^2}\,|u_n|,
\label{eq:psi1-resonant}
\end{equation}
where $Q_\psi = \Opsi_n/(2\gpsi_n)$ is the quality factor of the
$\psi$ mode.

\subsection{First-Order EM Correction: $(\mathbf{E}_1, \mathbf{B}_1)$}
\label{sec:EM1}

The modified EM fields at first order satisfy:
\begin{align}
\nabla^2\mathbf{E}_1 - \frac{n_0^2}{c^2}\ddot{\mathbf{E}}_1
&= \hat{\mathcal{F}}[\psi_1, \mathbf{E}_0, \mathbf{B}_0],
\label{eq:E1-equation}
\end{align}
where the source operator $\hat{\mathcal{F}}$ contains all terms
linear in $\psi_1$:
\begin{align}
\hat{\mathcal{F}} &= (1+\kap)\left[
\nabla(\nabla\psi_1\cdot\mathbf{E}_0)
- (\nabla\psi_1\cdot\nabla)\mathbf{E}_0\right] \notag\\
&\quad + (1-\kap)(\nabla\psi_1)\times(\nabla\times\mathbf{E}_0) \notag\\
&\quad + 2\frac{n_0}{c^2}\dot\psi_1\,\dot{\mathbf{E}}_0.
\label{eq:F-operator}
\end{align}

The key physical effects of $\psi_1$ on the EM fields are:
\begin{enumerate}[nosep]
\item \textbf{Frequency shift}: The resonant frequency shifts by
$\delta\omega/\omega = -\langle\psi_1\rangle_{\rm mode}$
\item \textbf{Mode mixing}: $\nabla\psi_1$ couples different
cavity modes
\item \textbf{Impedance change}: The cavity impedance shifts by
$\delta Z/Z = -\kap\langle\psi_1\rangle$
\end{enumerate}

For $|\psi_1| \sim 10^{-26}$, these corrections are:
\begin{equation}
\frac{|\mathbf{E}_1|}{|\mathbf{E}_0|} \sim |\psi_1|
\sim 10^{-26} \ll 1.
\label{eq:E1-smallness}
\end{equation}

\subsection{Second Order: $\psi_2$ (Back-Reaction Correction)}
\label{sec:psi2}

At second order, $\psi_2$ is sourced by the cross-terms between
zeroth and first-order fields:
\begin{align}
\nabla\cdot[\mufunction_0\,\nabla\psi_2]
+ \frac{1}{c_s^2}\mufunction_{t,0}\,\ddot\psi_2
&= -\frac{4\pi G\,\lam}{c^4}\left[
\delta u_{\rm EM}^{(1)} - \frac{\kap}{2}\delta\calB^{(1)}\right]
\notag\\
&\quad - \frac{4\pi G}{c^2}\,\psi_1\,
\frac{\partial\mufunction}{\partial\psi}\bigg|_0
\nabla^2\psi_0,
\label{eq:psi2-equation}
\end{align}
where:
\begin{align}
\delta u_{\rm EM}^{(1)} &= \epsilon_0 e^{(1+\kap)\psi_0}
\left[(1+\kap)\psi_1\,\frac{E_0^2}{2}
+ \mathbf{E}_0\cdot\mathbf{E}_1\right] \notag\\
&\quad + \frac{e^{-(1-\kap)\psi_0}}{2\mu_0}
\left[-(1-\kap)\psi_1\,B_0^2 + 2\mathbf{B}_0\cdot\mathbf{B}_1\right].
\label{eq:delta-uEM}
\end{align}

\paragraph{Magnitude estimate.} There are two contributions to $\psi_2$:
\begin{enumerate}[nosep]
\item Cross-term $\mathbf{E}_0\cdot\mathbf{E}_1$: This is
$\order{\varepsilon \times |\psi_1| \times u_{\rm EM}^{(0)}}$
$= \order{\varepsilon^2 \times u_{\rm EM}^{(0)}}$
\item Self-energy $\psi_1 E_0^2$: Also $\order{\varepsilon \times |\psi_1|
\times u_{\rm EM}^{(0)}} = \order{\varepsilon^2 \times u_{\rm EM}^{(0)}}$
\end{enumerate}

Therefore:
\begin{equation}
\boxed{
\frac{|\psi_2|}{|\psi_1|} \sim |\psi_1| + |\lam - 1|
\sim |\lam - 1| \lesssim 3 \times 10^{-5} \ll 1
}
\label{eq:psi2-ratio}
\end{equation}

The back-reaction correction is negligible compared to the first-order
perturbation, confirming the validity of the perturbative approach.

\subsection{Convergence of the Perturbation Series}\label{sec:convergence}

\begin{theorem}[Convergence of the EM--$\psi$ Perturbation Series]
\label{thm:convergence}
The perturbation series $\psi = \sum_{n=0}^\infty \psi_n$ converges
absolutely and uniformly for
\begin{equation}
|\lam - 1| < \varepsilon_{\rm max}
\equiv \min\left(\frac{c^4}{2G\,Q_\psi\,|G_n|\,|u_n|_{\rm max}},\;
\frac{2\gpsi_n}{\Opsi_n}\frac{\Mpsi_n\,\Opsi_n^2}{U_0\,H_n\,m}\right).
\label{eq:convergence-bound}
\end{equation}
\end{theorem}

\begin{proof}
Define the norm
$\|\psi_n\|_\infty = \sup_{\mathbf{x}} |\psi_n(\mathbf{x})|$.
At each order, the source for $\psi_{n+1}$ is bounded by:
\begin{equation}
\|{\rm source}_{n+1}\| \leq C_1\,|\lam-1|\,\|\psi_n\|_\infty\,
\bar{u}_{\rm EM} + C_2\,\|\psi_n\|_\infty^2\,\bar{u}_{\rm EM}/c^2,
\end{equation}
where $C_1$, $C_2$ are geometry-dependent constants of order unity.

The Green's function of the $\psi$ equation maps source to solution with
a gain factor $G_\psi$. The iterative estimate gives:
\begin{equation}
\|\psi_{n+1}\|_\infty \leq G_\psi\left[
C_1\,|\lam-1| + C_2\,\|\psi_n\|_\infty\right]
\|\psi_n\|_\infty\,\frac{\bar{u}_{\rm EM}}{c^2}.
\end{equation}

For $\|\psi_1\|_\infty \ll 1$ (which is guaranteed by the numerical
estimate~\eqref{eq:psi1-numerical}), the dominant term is the first,
giving a geometric series with ratio:
\begin{equation}
r = G_\psi\,C_1\,|\lam-1|\,\frac{\bar{u}_{\rm EM}}{c^2}.
\end{equation}

Convergence requires $r < 1$. Using
$G_\psi \sim Q_\psi/(\Mpsi_n\Opsi_n^2)$ and
$\bar{u}_{\rm EM} \sim U_0/V_{\rm shell}$:
\begin{equation}
r \sim \frac{|\lam-1|\,Q_\psi\,U_0}{\Mpsi_n\,\Opsi_n^2\,V_{\rm shell}\,c^2}
\sim |\lam-1|\,Q_\psi\,\frac{4\pi G\,U_0}{c^4\,V_{\rm shell}}
\,\frac{1}{\Opsi_n^2/c_s^2}.
\end{equation}

For $|\lam-1| \leq 3 \times 10^{-5}$, $Q_\psi \sim 500$,
$U_0 = 10^6\;\si{J}$, $V_{\rm shell} = 1\;\si{m^3}$:
\begin{equation}
r \sim 3\times10^{-5} \times 500 \times
\frac{4\pi \times 6.67\times10^{-11} \times 10^6}
{8.1\times10^{33} \times 1}
\sim 10^{-33} \ll 1.
\end{equation}

The series converges with overwhelming margin.
\end{proof}

\paragraph{Physical interpretation.} The convergence ratio $r \sim 10^{-33}$
reflects the enormous hierarchy between the EM energy scale and the
Planck energy scale. The gravitational coupling $G/c^4$ is so weak that
EM energy densities achievable in the laboratory produce $\psi$
perturbations that are fantastically small. The perturbation series is
not merely convergent; it is effectively truncated at first order for
all practical purposes.


%=============================================================================
%=============================================================================
\section{Stability Analysis}\label{sec:stability}
%=============================================================================
%=============================================================================

\subsection{Linearization Around the Self-Consistent Solution}

Let $(\psi_{\rm sc}, \mathbf{E}_{\rm sc}, \mathbf{B}_{\rm sc})$ denote
the self-consistent solution (to desired perturbative order). Perturb:
\begin{align}
\psi &= \psi_{\rm sc} + \delta\psi, \\
\mathbf{E} &= \mathbf{E}_{\rm sc} + \delta\mathbf{E}, \\
\mathbf{B} &= \mathbf{B}_{\rm sc} + \delta\mathbf{B}.
\end{align}

The linearized coupled system for
$(\delta\psi, \delta\mathbf{E}, \delta\mathbf{B})$ is:
\begin{align}
\hat{L}_\psi\,\delta\psi &= -\frac{4\pi G\,\lam}{c^4}
\left[\delta u_{\rm EM} - \frac{\kap}{2}\delta\calB\right],
\label{eq:lin-psi}\\
\hat{L}_{\rm EM}\,\delta\mathbf{E} &= \hat{\mathcal{F}}
[\delta\psi, \mathbf{E}_{\rm sc}, \mathbf{B}_{\rm sc}],
\label{eq:lin-EM}
\end{align}
where $\hat{L}_\psi$ is the linearized $\psi$ wave operator and
$\hat{L}_{\rm EM}$ is the linearized EM wave operator, both evaluated
on the self-consistent background.

\subsection{Modal Decomposition of the Coupled System}

Expand perturbations in the joint $\psi$-EM eigenbasis:
\begin{align}
\delta\psi &= \sum_n a_n(t)\,u_n(\mathbf{x}), \\
\delta\mathbf{E} &= \sum_\alpha b_\alpha(t)\,\mathbf{e}_\alpha(\mathbf{x}),
\end{align}
where $\{u_n\}$ are the $\psi$-mode eigenfunctions and
$\{\mathbf{e}_\alpha\}$ are the EM cavity modes.

Projecting onto modes, the coupled amplitude equations become:
\begin{align}
\ddot{a}_n + 2\gpsi_n\dot{a}_n + \Opsi_n^2\,a_n
&= \sum_\alpha C_{n\alpha}\,b_\alpha
+ \sum_\alpha D_{n\alpha}(t)\,a_n,
\label{eq:coupled-an}\\
\ddot{b}_\alpha + 2\gamma_\alpha\dot{b}_\alpha + \omega_\alpha^2\,b_\alpha
&= \sum_n F_{\alpha n}\,a_n,
\label{eq:coupled-balpha}
\end{align}
where the coupling matrices are:
\begin{align}
C_{n\alpha} &= \frac{\lam}{c^2\Mpsi_n}\int u_n(\mathbf{x})\left[
\epsilon_{\rm sc}\,\mathbf{E}_{\rm sc}\cdot\mathbf{e}_\alpha
+ \frac{\mathbf{B}_{\rm sc}\cdot\mathbf{b}_\alpha}{\mu_{\rm m,sc}}
\right]d^3x, \label{eq:C-matrix}\\
F_{\alpha n} &= \int \mathbf{e}_\alpha^*\cdot
\hat{\mathcal{F}}[u_n, \mathbf{E}_{\rm sc}, \mathbf{B}_{\rm sc}]
\,d^3x. \label{eq:F-matrix}
\end{align}

\subsection{Eigenfrequency Analysis}\label{sec:eigenfreq}

The coupled system~\eqref{eq:coupled-an}--\eqref{eq:coupled-balpha}
has eigenfrequencies determined by the secular equation:
\begin{equation}
\det\begin{pmatrix}
-\omega^2 + 2i\gpsi_n\omega + \Opsi_n^2 - D_{nn}
& -C_{n\alpha} \\
-F_{\alpha n}
& -\omega^2 + 2i\gamma_\alpha\omega + \omega_\alpha^2
\end{pmatrix} = 0.
\label{eq:secular}
\end{equation}

For a single $\psi$ mode coupled to a single EM mode, this reduces to:
\begin{equation}
(\omega^2 - \Opsi_n^2 - 2i\gpsi_n\omega + D_{nn})
(\omega^2 - \omega_\alpha^2 - 2i\gamma_\alpha\omega)
= C_{n\alpha}F_{\alpha n}.
\label{eq:secular-2mode}
\end{equation}

\paragraph{Weak coupling limit.} When $|C_{n\alpha}F_{\alpha n}|
\ll \Opsi_n^2\omega_\alpha^2$ (which is guaranteed by the $G/c^4$
suppression), the eigenfrequencies split into:
\begin{align}
\omega_\pm^{(\psi)} &\simeq \pm\Opsi_n + i\gpsi_n
+ \frac{C_{n\alpha}F_{\alpha n}}
{2\Opsi_n(\Opsi_n^2 - \omega_\alpha^2 + 2i\gamma_\alpha\Opsi_n)},
\label{eq:omega-psi-shifted}\\
\omega_\pm^{(\rm EM)} &\simeq \pm\omega_\alpha + i\gamma_\alpha
+ \frac{C_{n\alpha}F_{\alpha n}}
{2\omega_\alpha(\omega_\alpha^2 - \Opsi_n^2 + 2i\gpsi_n\omega_\alpha)}.
\label{eq:omega-EM-shifted}
\end{align}

\begin{proposition}[Stability Below Parametric Threshold]
\label{prop:stability}
For $|\lam - 1| < |\lam - 1|_{\rm min}$ (below the parametric
instability threshold), all eigenfrequencies have
$\mathrm{Im}(\omega) > 0$ (using the $e^{+i\omega t}$ convention),
corresponding to exponential decay. The system is linearly stable.
\end{proposition}

\begin{proof}
The imaginary parts of the shifted eigenfrequencies are:
\begin{align}
\mathrm{Im}(\omega_\pm^{(\psi)}) &= \gpsi_n
- \frac{|C_{n\alpha}F_{\alpha n}|\,\gamma_\alpha}
{|\Opsi_n^2 - \omega_\alpha^2|^2 + 4\gamma_\alpha^2\Opsi_n^2}
+ \order{\varepsilon^4},
\end{align}
and similarly for the EM mode. The correction to the damping rate is:
\begin{equation}
|\delta\gpsi_n| \sim \frac{|C_{n\alpha}||F_{\alpha n}|\gamma_\alpha}
{|\Opsi_n^2 - \omega_\alpha^2|^2}
\sim \frac{(\lam-1)^2\,u_{\rm EM}^2\,\gamma_\alpha}
{c^4\,\Mpsi_n^2\,|\Opsi_n^2 - \omega_\alpha^2|^2}.
\end{equation}

This is of order $\varepsilon^2 \times (G/c^4)^2 \times u_{\rm EM}^2$,
which is negligible compared to $\gpsi_n$. Therefore
$\mathrm{Im}(\omega) > 0$ for all modes, and the system is stable.

The exception is the parametric channel: when
$|\lam-1| > |\lam-1|_{\rm min}$, the time-dependent coupling
$D_{nn}(t)$ drives parametric instability through the Mathieu mechanism.
Below threshold, the time-averaged $D_{nn}$ produces only a small
frequency shift, not instability.
\end{proof}

\subsection{What Happens Above Threshold?}\label{sec:above-threshold}

If $|\lam - 1| > |\lam - 1|_{\rm min}$, the parametric instability
grows with rate:
\begin{equation}
\Gamma = \frac{1}{2}h\,\Opsi_n - \gpsi_n,
\label{eq:instability-rate}
\end{equation}
where $h = |\lam-1|\,U_0\,H_n\,m/(\Mpsi_n\,\Opsi_n^2)$.

The instability is self-limiting because as $\psi_1$ grows:
\begin{enumerate}[nosep]
\item The resonant frequency shifts:
$\Opsi_n \to \Opsi_n + \delta\Opsi \propto \psi_1^2$, detuning the
parametric resonance
\item Nonlinear mode-mode coupling transfers energy to off-resonant
modes, increasing effective damping
\item The modified $\epsilon(\psi)$, $\mu_{\rm m}(\psi)$ shift the EM
cavity frequency, detuning the EM drive
\end{enumerate}

The saturation amplitude is determined by the detuning condition:
\begin{equation}
|\psi_1|_{\rm sat} \sim \frac{\gpsi_n}{\Opsi_n}
\frac{1}{|\partial\Opsi_n/\partial\psi_1|}
\sim \frac{\gpsi_n}{\Opsi_n}.
\label{eq:saturation}
\end{equation}

For $\gpsi_n/\Opsi_n \sim 10^{-3}$, the saturation amplitude is
$|\psi_1|_{\rm sat} \sim 10^{-3}$, corresponding to a local
acceleration perturbation of
$\delta a \sim c^2\psi_1/(2L) \sim 10^{-3} \times c^2/(2L)$.

\paragraph{No thermal runaway.} The system cannot run away to
arbitrarily large $\psi$ values because:
\begin{enumerate}[nosep]
\item The convexity of $\Wfunc$ ensures the energy functional is
bounded below
\item Nonlinear detuning provides negative feedback
\item Energy conservation (Section~\ref{sec:energy}) bounds the total
available energy
\end{enumerate}

\begin{tcolorbox}[colback=red!5!white, colframe=red!60!black,
title=Maximum Stable Perturbation]
The maximum $\psi$ perturbation before the system departs the linear
regime is:
\begin{equation}
|\psi|_{\rm max,linear} \sim \frac{\gpsi_n}{\Opsi_n}
\sim 10^{-3}.
\end{equation}
Beyond this, nonlinear effects dominate but the system saturates
rather than running away, due to the convexity of $\Wfunc$ and
energy conservation. The maximum attainable perturbation is bounded
by the total EM energy available:
\begin{equation}
|\psi|_{\rm max} \lesssim \frac{4\pi G\,U_{\rm total}}{c^4\,R}
\sim 10^{-22}\;\text{for laboratory scales}.
\end{equation}
\end{tcolorbox}


%=============================================================================
%=============================================================================
\section{Energy Conservation}\label{sec:energy}
%=============================================================================
%=============================================================================

\subsection{Derivation from the Action via Noether's Theorem}

The total action $S = S_\psi + S_{\rm EM} + S_{\rm matter}$ is invariant
under time translations $t \to t + \delta t$ (assuming no explicit time
dependence in the background). By Noether's theorem, the conserved
energy is:
\begin{equation}
\boxed{
E_{\rm total} = E_\psi + E_{\rm EM} + E_{\rm matter} + E_{\rm int}
}
\label{eq:total-energy}
\end{equation}
where each component is identified below.

\subsubsection{$\psi$-Field Energy}

From the stress-energy tensor of $S_\psi$:
\begin{equation}
E_\psi = \int d^3x\left[
\frac{\astar^2}{8\pi G}\Wfunc'(y)\frac{|\nabla\psi|^2}{\astar^2}
- \frac{\astar^2}{8\pi G}\Wfunc(y)
+ \frac{1}{2c_s^2}\mufunction_t\,\dot\psi^2
\right],
\label{eq:E-psi}
\end{equation}
where $y = |\nabla\psi|^2/\astar^2$. In the Newtonian limit
($\Wfunc(y) \approx y$):
\begin{equation}
E_\psi \to \int d^3x\;\frac{|\nabla\psi|^2}{8\pi G}
= -\frac{1}{2}\int d^3x\;\rho\,\Phi,
\label{eq:E-psi-Newton}
\end{equation}
recovering the standard gravitational self-energy.

\subsubsection{EM Field Energy}

\begin{equation}
E_{\rm EM} = \int d^3x\left[
\frac{\epsilon(\psi)}{2}E^2 + \frac{B^2}{2\mu_{\rm m}(\psi)}
\right]
= \int d^3x\;u_{\rm EM}.
\label{eq:E-EM}
\end{equation}

Note: $E_{\rm EM}$ depends on $\psi$ through the constitutive relations.
This is the source of the interaction energy.

\subsubsection{Matter Kinetic Energy}

\begin{equation}
E_{\rm matter} = \sum_a \frac{m_a c^2}{\sqrt{1-v_a^2/c^2}}\,
e^{-\psi(\mathbf{x}_a)/2}
\approx \sum_a \left(m_a c^2 + \frac{1}{2}m_a v_a^2
- \frac{1}{2}m_a c^2\psi_a\right).
\label{eq:E-matter}
\end{equation}

\subsubsection{Interaction Energy}

The interaction energy arises from the $\psi$-dependence of the
matter and EM energies:
\begin{equation}
E_{\rm int} = -\frac{c^2}{2}\int d^3x\;\psi\,\rho
- \frac{\lam}{c^2}\int d^3x\;\psi\left(
u_{\rm EM} - \frac{\kap}{2}\calB\right).
\label{eq:E-int}
\end{equation}

\subsection{The Modified Poynting Theorem}\label{sec:poynting}

The energy continuity equations for each sector are:

\paragraph{EM sector:}
\begin{equation}
\frac{\partial u_{\rm EM}}{\partial t}
+ \nabla\cdot(\mathbf{E}\times\mathbf{H})
= -\mathbf{J}\cdot\mathbf{E}
- \frac{\kap}{2}\dot\psi\,\calB
- \dot\psi\,u_{\rm EM}.
\label{eq:poynting-EM}
\end{equation}

The last two terms represent energy exchange between EM and $\psi$:
\begin{itemize}[nosep]
\item $-\dot\psi\,u_{\rm EM}$: energy flow due to changing refractive
index (both $\epsilon$ and $\mu_{\rm m}$ change)
\item $-(\kap/2)\dot\psi\,\calB$: differential $E$-vs-$B$ energy
exchange through the dual-sector coupling
\end{itemize}

\paragraph{$\psi$ sector:}
\begin{equation}
\frac{\partial u_\psi}{\partial t}
+ \nabla\cdot\mathbf{S}_\psi
= -\frac{c^2}{2}\dot\psi\,\rho
+ \frac{\lam}{c^2}\dot\psi\left(u_{\rm EM} - \frac{\kap}{2}\calB\right),
\label{eq:poynting-psi}
\end{equation}
where the $\psi$-field Poynting vector is:
\begin{equation}
\mathbf{S}_\psi = -\frac{c^2}{4\pi G}\mufunction\,\dot\psi\,\nabla\psi.
\label{eq:S-psi-poynting}
\end{equation}

\paragraph{Combined:} Adding Eqs.~\eqref{eq:poynting-EM}
and~\eqref{eq:poynting-psi} with the matter kinetic energy rate:
\begin{equation}
\frac{\partial}{\partial t}(u_\psi + u_{\rm EM} + u_{\rm kin})
+ \nabla\cdot(\mathbf{S}_\psi + \mathbf{S}_{\rm EM})
= -\mathbf{J}\cdot\mathbf{E}
+ (1-\lam)\frac{\dot\psi}{c^2}\left(u_{\rm EM}
- \frac{\kap}{2}\calB\right).
\label{eq:combined-energy}
\end{equation}

For $\lam = 1$: the exchange term vanishes, and energy is conserved
in the standard way. For $\lam \neq 1$: the exchange term represents
energy flow between the EM field and the $\psi$ field, mediated by the
anomalous coupling. But total energy is still conserved:

\begin{proposition}[Total Energy Conservation]
The total energy $E_{\rm total}$ defined by
Eq.~\eqref{eq:total-energy} is exactly conserved:
$dE_{\rm total}/dt = 0$.
\end{proposition}

\begin{proof}
This follows directly from time-translation invariance of the action
$S$ via Noether's theorem. Alternatively, one can verify by direct
computation that the exchange terms between sectors sum to zero
identically. The $\psi$ sector gains energy at rate
$\lam\dot\psi(u_{\rm EM} - \kap\calB/2)/c^2$, while the EM sector
loses energy at the same rate. The matter sector exchanges energy
with $\psi$ at rate $-(c^2/2)\dot\psi\rho$. All exchange terms are
antisymmetric (what one sector gains, another loses), ensuring exact
conservation.
\end{proof}

\subsection{The Energy Budget for the Psi-Bubble Device}
\label{sec:energy-budget}

\begin{tcolorbox}[colback=green!5!white, colframe=green!60!black,
title=Critical Question: Where Does the Kinetic Energy Come From?]
If the vehicle accelerates to velocity $v$ and gains kinetic energy
$\frac{1}{2}Mv^2$, what loses an equal amount of energy?
\end{tcolorbox}

The answer emerges from the energy conservation
equation~\eqref{eq:combined-energy}. The energy flow path is:

\begin{enumerate}
\item \textbf{EM power supply} provides energy $P_{\rm in}$ to maintain
the cavity fields against losses.

\item \textbf{EM-to-$\psi$ conversion}: A fraction of the EM energy
converts to $\psi$-field energy through the $(\lam-1)$ coupling:
\begin{equation}
P_{\rm EM\to\psi} = \frac{(\lam-1)}{c^2}\int \dot\psi
\left(u_{\rm EM} - \frac{\kap}{2}\calB\right)d^3x.
\label{eq:P-EM-psi}
\end{equation}

\item \textbf{$\psi$-gradient creation}: The EM-sourced $\psi$
perturbation creates a local $\psi$ gradient, which produces an
acceleration field $\avec = (c^2/2)\nabla\psi$.

\item \textbf{Vehicle acceleration}: The vehicle (shell + payload)
accelerates in the $\psi$ gradient it has created, gaining kinetic
energy at rate:
\begin{equation}
\dot{E}_{\rm kin} = M\,\avec\cdot\mathbf{v}
= \frac{Mc^2}{2}(\nabla\psi)\cdot\mathbf{v}.
\label{eq:Ekin-rate}
\end{equation}

\item \textbf{Recoil in $\psi$ field}: The $\psi$ field itself carries
energy away as scalar radiation:
\begin{equation}
P_{\rm rad}^\psi = \oint \mathbf{S}_\psi\cdot d\mathbf{A}
= -\frac{c^2}{4\pi G}\oint \mufunction\,\dot\psi\,
(\nabla\psi)\cdot d\mathbf{A}.
\label{eq:P-rad}
\end{equation}
\end{enumerate}

The complete energy budget at steady state is:
\begin{equation}
\boxed{
P_{\rm in} = P_{\rm cavity\;loss} + P_{\rm EM\to\psi}
= P_{\rm cavity\;loss} + \dot{E}_{\rm kin} + P_{\rm rad}^\psi
+ P_{\psi\;dissip}
}
\label{eq:energy-budget}
\end{equation}

\paragraph{The energy comes from the EM field.} This is the definitive
answer to the energy question. The kinetic energy of the vehicle comes
from the electromagnetic energy stored in the cavity, mediated by the
$\psi$-field coupling. The device is fundamentally an EM-to-kinetic
energy converter, not a ``reactionless drive'' or perpetual motion
machine. The power supply must continuously provide energy to:
\begin{enumerate}[nosep]
\item Maintain the cavity fields against ohmic/radiation losses
\item Replace the energy converted to $\psi$-field energy
\item The $\psi$-field energy is then distributed between kinetic
energy of the vehicle and radiated $\psi$ waves
\end{enumerate}

\paragraph{Efficiency.} The propulsive efficiency is:
\begin{equation}
\eta_{\rm prop} = \frac{\dot{E}_{\rm kin}}{P_{\rm in}}
= \frac{P_{\rm EM\to\psi} - P_{\rm rad}^\psi - P_{\psi\;dissip}}
{P_{\rm cavity\;loss} + P_{\rm EM\to\psi}}.
\label{eq:efficiency}
\end{equation}

For $P_{\rm cavity\;loss} \gg P_{\rm EM\to\psi}$ (which is the case
for $|\lam-1| \ll 1$):
\begin{equation}
\eta_{\rm prop} \sim \frac{|\lam-1|\,U_0\,Q_\psi}
{P_{\rm cavity\;loss}\,c^2/G}
\sim |\lam-1| \times 10^{-20},
\label{eq:efficiency-estimate}
\end{equation}
which is extremely small, consistent with the enormous ratio between
electromagnetic and gravitational energy scales.


%=============================================================================
%=============================================================================
\section{Momentum Conservation}\label{sec:momentum}
%=============================================================================
%=============================================================================

\subsection{The Momentum Budget}

\begin{tcolorbox}[colback=green!5!white, colframe=green!60!black,
title=Critical Question: What Carries the Recoil Momentum?]
If the vehicle gains momentum $M\mathbf{v}$, what carries the equal
and opposite momentum $-M\mathbf{v}$?
\end{tcolorbox}

\subsection{Momentum Density in Each Sector}

\subsubsection{EM Momentum}

The EM momentum density is the standard Poynting vector divided by $c^2$:
\begin{equation}
\mathbf{g}_{\rm EM} = \frac{\mathbf{E}\times\mathbf{H}}{c^2}
= \frac{\mathbf{E}\times\mathbf{B}}{c^2\,\mu_{\rm m}(\psi)}.
\label{eq:g-EM}
\end{equation}

For a standing-wave cavity, the time-averaged EM momentum vanishes:
$\langle\mathbf{g}_{\rm EM}\rangle = 0$. For a rotating mode, it does
not vanish but has azimuthal (angular) character, not linear.

\subsubsection{$\psi$-Field Momentum}

The $\psi$-field momentum density follows from the $T^{0i}$ component
of the stress-energy tensor:
\begin{equation}
\mathbf{g}_\psi = -\frac{1}{4\pi G}\mufunction\,\dot\psi\,\nabla\psi.
\label{eq:g-psi}
\end{equation}

This is the key object: the $\psi$ field carries momentum whenever
it is both time-varying ($\dot\psi \neq 0$) and spatially varying
($\nabla\psi \neq 0$).

\paragraph{Physical interpretation.} In GR, gravitational waves carry
momentum through the Isaacson stress-energy tensor. In DFD, the $\psi$
field carries momentum through~\eqref{eq:g-psi}. The mathematical
structure is identical to the momentum density of a massive scalar field.

\subsubsection{Matter Momentum}

\begin{equation}
\mathbf{P}_{\rm matter} = \sum_a m_a\,e^{-\psi_a/2}\,
\frac{\mathbf{v}_a}{\sqrt{1-v_a^2/c^2}}
\approx M\mathbf{v}_{\rm vehicle}.
\label{eq:P-matter}
\end{equation}

\subsection{Total Momentum Conservation}

\begin{proposition}[Total Momentum Conservation]
The total momentum
\begin{equation}
\mathbf{P}_{\rm total} = \int d^3x\,(\mathbf{g}_{\rm EM}
+ \mathbf{g}_\psi) + \mathbf{P}_{\rm matter}
\label{eq:P-total}
\end{equation}
is exactly conserved: $d\mathbf{P}_{\rm total}/dt = 0$.
\end{proposition}

\begin{proof}
The action $S$ is invariant under spatial translations
$\mathbf{x} \to \mathbf{x} + \delta\mathbf{x}$ (for fields on flat
$\mathbb{R}^3$). By Noether's theorem, the associated conserved current
yields momentum conservation. The momentum continuity equation is:
\begin{equation}
\frac{\partial g_i}{\partial t} + \partial_j T_{ij} = 0,
\end{equation}
where $T_{ij}$ is the total spatial stress tensor. Integrating over all
space (with appropriate fall-off conditions at infinity):
\begin{equation}
\frac{d}{dt}\int d^3x\,g_i = -\oint T_{ij}\,dA_j \to 0
\quad\text{(for localized sources)}.
\end{equation}
\end{proof}

\subsection{Momentum Flow in the Psi-Bubble Device}

When the device creates an asymmetric $\psi$ perturbation and the
vehicle accelerates:
\begin{equation}
\frac{d\mathbf{P}_{\rm matter}}{dt} = M\frac{c^2}{2}\nabla\psi
= -\frac{d}{dt}\int d^3x\;(\mathbf{g}_\psi + \mathbf{g}_{\rm EM}).
\label{eq:momentum-balance}
\end{equation}

The recoil momentum is carried by the $\psi$ field. Specifically, as
the device sources an asymmetric $\psi$ perturbation (via the asymmetric
EM energy distribution), the $\psi$ field develops a time-varying,
spatially varying pattern that carries momentum flux:
\begin{equation}
\mathbf{F}_{\rm recoil} = -\oint T_{ij}^\psi\,dA_j
= \oint \frac{c^4}{4\pi G}\mufunction\left[
\frac{1}{2}|\nabla\psi|^2\hat{n}_i - (\nabla\psi)_i(\nabla\psi)_j\hat{n}_j
\right]dA.
\label{eq:psi-momentum-flux}
\end{equation}

In the radiation zone (far from the device), this becomes the
$\psi$-wave momentum flux:
\begin{equation}
\mathbf{F}_{\rm rad}^\psi = \frac{P_{\rm rad}^\psi}{c_s}\,\hat{\mathbf{n}},
\label{eq:rad-momentum}
\end{equation}
where $P_{\rm rad}^\psi$ is the radiated $\psi$-wave power and $c_s$
is the propagation speed. The momentum carried by $\psi$ radiation
provides the recoil: $\mathbf{P}_{\rm vehicle} = -\mathbf{P}_{\rm rad}^\psi$.

\paragraph{Analogy with photon propulsion.} The device is analogous to
a photon rocket, but emitting $\psi$ waves instead of photons. For a
photon rocket, $\mathbf{F} = P/c$; here,
$\mathbf{F} = P_{\rm rad}^\psi/c_s$.


%=============================================================================
%=============================================================================
\section{Radiation of $\psi$ Waves}\label{sec:psi-radiation}
%=============================================================================
%=============================================================================

\subsection{The $\psi$-Wave Equation in the Radiation Zone}

In the far zone ($r \gg R_{\rm device}$), the $\psi$-wave equation
linearizes to:
\begin{equation}
\nabla^2\psi - \frac{1}{c_s^2}\ddot\psi
= -\frac{4\pi G}{c^2}\,\rho_{\rm eff}(\mathbf{x},t),
\label{eq:psi-wave-eq}
\end{equation}
where $\rho_{\rm eff}$ includes both matter and (if $\lam \neq 1$)
EM energy density sources. The retarded Green's function solution is:
\begin{equation}
\psi(\mathbf{x},t) = \frac{G}{c^2}\int
\frac{\rho_{\rm eff}(\mathbf{x}',t-|\mathbf{x}-\mathbf{x}'|/c_s)}
{|\mathbf{x}-\mathbf{x}'|}\,d^3x'.
\label{eq:retarded-solution}
\end{equation}

\subsection{$\psi$-Wave Luminosity}\label{sec:luminosity}

For a source with time-varying quadrupole moment
$Q_{ij}(t) = \int \rho_{\rm eff}\,x_i x_j\,d^3x$, the radiated power
in $\psi$ waves is:

\paragraph{Monopole radiation (scalar specific).}
Unlike the tensor GW sector (which requires quadrupole variation), the
scalar $\psi$ field can radiate at monopole order. The monopole moment is:
\begin{equation}
\mathcal{M}(t) = \int \rho_{\rm eff}(\mathbf{x},t)\,d^3x.
\label{eq:monopole}
\end{equation}

For the psi-bubble device, the EM energy density oscillates at $2\omega$,
giving:
\begin{equation}
\dot{\mathcal{M}} = \frac{\lam}{c^2}\int \dot{u}_{\rm EM}\,d^3x
= \frac{\lam}{c^2}\,2\omega\,\hat{U}\sin(2\omega t),
\label{eq:Mdot}
\end{equation}
where $\hat{U}$ is the amplitude of the oscillating EM energy component.

The monopole $\psi$-wave luminosity is:
\begin{equation}
\boxed{
L_\psi^{(\rm mono)} = \frac{G}{c_s^3}\langle\ddot{\mathcal{M}}^2\rangle
= \frac{2G\,\lam^2\,\omega^2\,\hat{U}^2}{c^4\,c_s^3}
}
\label{eq:L-psi-mono}
\end{equation}

\paragraph{Numerical estimate.} For $\omega/(2\pi) = 1.3\;\si{GHz}$,
$\hat{U} = mU_0 = 0.1 \times 10^6 = 10^5\;\si{J}$, $c_s = c$:
\begin{align}
L_\psi^{(\rm mono)} &= \frac{2 \times 6.67\times10^{-11}
\times (2\pi\times1.3\times10^9)^2 \times (10^5)^2}
{(3\times10^8)^4 \times (3\times10^8)^3} \notag\\
&\sim \frac{2 \times 6.67\times10^{-11} \times 6.7\times10^{19}
\times 10^{10}}{8.1\times10^{33} \times 2.7\times10^{25}} \notag\\
&\sim \frac{8.9\times10^{19}}{2.2\times10^{59}} \notag\\
&\sim 4 \times 10^{-40}\;\si{W}.
\label{eq:L-psi-numerical}
\end{align}

This is extraordinarily small---undetectable by any conceivable
instrument.

\subsection{Frequency and Spectrum}

The dominant $\psi$-wave frequency is:
\begin{equation}
f_\psi = 2f_{\rm EM} = 2 \times 1.3\;\si{GHz} = 2.6\;\si{GHz},
\label{eq:f-psi}
\end{equation}
because the EM energy density oscillates at twice the EM field frequency
(being quadratic in the fields).

If the device operates at variable power, the $\psi$-wave spectrum
includes:
\begin{itemize}[nosep]
\item $f = 2f_{\rm EM}$: fundamental (from $E^2$, $B^2$ oscillation)
\item $f = 0$: DC component (from time-averaged energy density)
\item $f = 2f_{\rm EM} \pm f_{\rm mod}$: sidebands from amplitude
modulation
\end{itemize}

\subsection{Braking Force from $\psi$ Radiation}

The radiation reaction force on the device is:
\begin{equation}
\mathbf{F}_{\rm rad} = -\frac{L_\psi}{c_s}\,\hat{\mathbf{v}}
\sim -4\times10^{-40}/c \sim 10^{-48}\;\si{N}.
\label{eq:braking-force}
\end{equation}

This is negligible compared to any other force in the system. The
$\psi$-wave radiation does not provide a significant braking mechanism
for any conceivable laboratory device.

\subsection{Detectability}\label{sec:detectability}

The $\psi$-wave strain at distance $r$ from the device is:
\begin{equation}
|\psi_{\rm wave}(r)| = \frac{G\,\lam\,\hat{U}}{c^4\,r}
\sim \frac{6.67\times10^{-11} \times 10^5}
{8.1\times10^{33} \times r}
\sim \frac{10^{-39}}{r},
\label{eq:psi-strain}
\end{equation}
which at $r = 1\;\si{m}$ gives $|\psi_{\rm wave}| \sim 10^{-39}$.
This is far below the projected sensitivity of any $\psi$-field
detector ($\sim 10^{-18}$ for optimized configurations), making the
$\psi$ radiation from a single device undetectable.


%=============================================================================
%=============================================================================
\section{Thermodynamic Consistency}\label{sec:thermodynamics}
%=============================================================================
%=============================================================================

\subsection{First Law of Thermodynamics}

The first law requires that the total internal energy change equals
heat added minus work done:
\begin{equation}
dE = \delta Q - \delta W.
\label{eq:first-law}
\end{equation}

For the psi-bubble device:
\begin{itemize}[nosep]
\item $dE = d(E_\psi + E_{\rm EM} + E_{\rm kin})$
\item $\delta Q = P_{\rm in}\,dt$ (electrical power input)
\item $\delta W = F_{\rm thrust}\,ds$ (work done on payload)
\end{itemize}

Energy conservation (Eq.~\eqref{eq:energy-budget}) ensures the first
law is satisfied identically. The device converts EM energy to kinetic
energy (and waste heat) through the $\psi$-field intermediary. No energy
is created or destroyed.

\subsection{Second Law of Thermodynamics}

The second law requires $\Delta S_{\rm total} \geq 0$ for any process.

\paragraph{Entropy sources in the device:}
\begin{enumerate}[nosep]
\item \textbf{Cavity ohmic losses}: $\dot{S}_{\rm ohmic}
= P_{\rm loss}/T_{\rm cavity}$
\item \textbf{$\psi$-field damping}: $\dot{S}_{\psi\,{\rm damp}}
= P_{\psi\,{\rm dissip}}/T_{\rm eff}$
\item \textbf{$\psi$-wave radiation}: $\dot{S}_{\rm rad}
= P_{\rm rad}^\psi/T_{\rm eff}^\psi$
(entropy carried by radiation)
\end{enumerate}

All three contributions are non-negative, so:
\begin{equation}
\dot{S}_{\rm total} = \dot{S}_{\rm ohmic} + \dot{S}_{\psi\,{\rm damp}}
+ \dot{S}_{\rm rad} \geq 0.
\label{eq:entropy-production}
\end{equation}

The second law is satisfied with strict inequality (the process is
irreversible due to ohmic losses and radiation).

\subsection{No Violation of the Third Law}

The device does not create a zero-entropy state or extract work from
a system at absolute zero. The superconducting cavities operate at
$T \approx 2\;\si{K}$, well above absolute zero.

\subsection{Maximum Thermodynamic Efficiency}

The Carnot-like maximum efficiency for the EM-to-kinetic conversion is:
\begin{equation}
\eta_{\rm max} = \frac{P_{\rm EM\to\psi}}{P_{\rm in}}
= \frac{(\lam-1)^2\,U_0^2\,H^2\,m^2\,\Opsi^2}
{4\Mpsi\gpsi\,P_{\rm in}},
\label{eq:eta-max}
\end{equation}
which is the ratio of the EM-to-$\psi$ coupling power to the total
input power. For $|\lam-1| \sim 10^{-5}$:
\begin{equation}
\eta_{\rm max} \sim 10^{-10} \times \frac{U_0}{P_{\rm in}/\Opsi}
\sim 10^{-10} \times Q_{\rm cav} \times \frac{U_0\,\Opsi}{P_{\rm rf}},
\end{equation}
which is very small for laboratory parameters.

\paragraph{Note.} This extreme inefficiency reflects the weakness of
the gravitational coupling, not a fundamental thermodynamic limit. In
principle, with $|\lam-1|$ close to unity (which is excluded by current
bounds), the efficiency could approach unity. The device is
thermodynamically legal but gravitationally feeble.

\subsection{Comparison with Other Propulsion Concepts}

\begin{table}[h]
\centering
\caption{Thermodynamic comparison of propulsion concepts.}
\label{tab:thermo-comparison}
\begin{ruledtabular}
\begin{tabular}{lcc}
Concept & Conserves $E$? & Conserves $\mathbf{P}$? \\
\hline
Chemical rocket & Yes & Yes (exhaust) \\
Ion drive & Yes & Yes (exhaust) \\
Solar sail & Yes & Yes (photon recoil) \\
Photon rocket & Yes & Yes (photon momentum) \\
EmDrive (claimed) & ??? & \textbf{No} \\
Psi-bubble (DFD) & \textbf{Yes} & \textbf{Yes} ($\psi$ radiation) \\
\end{tabular}
\end{ruledtabular}
\end{table}

The psi-bubble device conserves both energy and momentum. The ``exhaust''
is $\psi$-wave radiation, which carries both energy and momentum away
from the device. The device is analogous to a photon rocket, but with
$\psi$ waves replacing photons as the propellant. The thrust-to-power
ratio is:
\begin{equation}
\frac{F}{P} = \frac{1}{c_s} \times \eta_{\rm prop}
\sim \frac{10^{-10}}{c} \sim 10^{-18}\;\si{N/W},
\label{eq:F-over-P}
\end{equation}
compared to $2/c \sim 7 \times 10^{-9}\;\si{N/W}$ for a perfect
photon rocket.


%=============================================================================
%=============================================================================
\section{Summary and Key Results}\label{sec:summary}
%=============================================================================
%=============================================================================

We have solved the fully coupled DFD + Maxwell system for the
psi-bubble device. The principal results are:

\begin{enumerate}
\item \textbf{Complete coupled system} (Section~\ref{sec:coupled-system}):
The exact field equations are given by
Eqs.~\eqref{eq:psi-field-exact},~\eqref{eq:gauss}--\eqref{eq:faraday},
and~\eqref{eq:matter-eom}, with $\psi$-dependent constitutive
relations~\eqref{eq:eps}--\eqref{eq:mum}.

\item \textbf{Perturbative solution converges}
(Section~\ref{sec:perturbative}): The perturbation series has
convergence ratio $r \sim 10^{-33}$, far below unity. For
$|\lam-1| \leq 3\times10^{-5}$, first-order perturbation theory is
exact to $\sim 33$ decimal places.

\item \textbf{$\psi_1$ is computable}
(Section~\ref{sec:psi1}): The DC component is
$\psi_1^{(\rm DC)} \sim (\lam-1) \times 10^{-21}$ at device
surface; the resonant oscillating component is enhanced by $Q_\psi$.

\item \textbf{$|\psi_2/\psi_1| \sim |\lam-1| \ll 1$}
(Section~\ref{sec:psi2}): The back-reaction correction is negligible.

\item \textbf{System is stable below parametric threshold}
(Section~\ref{sec:stability}): All coupled eigenfrequencies have
positive imaginary parts (damped). Above threshold, nonlinear saturation
prevents runaway.

\item \textbf{Energy is exactly conserved}
(Section~\ref{sec:energy}): Total energy
$E_\psi + E_{\rm EM} + E_{\rm matter} + E_{\rm int} = {\rm const}$.
The kinetic energy of the vehicle comes from the EM field, mediated
by $\psi$-coupling. There is no violation of energy conservation.

\item \textbf{Momentum is exactly conserved}
(Section~\ref{sec:momentum}): The recoil momentum is carried by the
$\psi$ field (both near-field and radiated components). The device is
analogous to a photon rocket emitting $\psi$ waves.

\item \textbf{$\psi$-wave luminosity is negligible}
(Section~\ref{sec:psi-radiation}): $L_\psi \sim 10^{-40}\;\si{W}$
for laboratory-scale devices. The radiation is undetectable and the
braking force is negligible.

\item \textbf{No thermodynamic violations}
(Section~\ref{sec:thermodynamics}): All three laws are satisfied.
The maximum efficiency is $\sim |\lam-1|^2 \times (G\,U_0/c^4)$,
which is extremely small, reflecting the weakness of gravity.
\end{enumerate}

\paragraph{The bottom line for skeptics.} The psi-bubble device, within
the DFD framework with $|\lam-1| \lesssim 3\times10^{-5}$, is a
thermodynamically consistent, energy-conserving, momentum-conserving
machine. It is not a perpetual motion device. It is not reactionless.
It converts EM energy to kinetic energy with an efficiency suppressed
by $(G/c^4)$, making the effect fantastically small for current
technology. The ``propulsion'' aspect is real in principle but
unmeasurable in practice unless $|\lam-1|$ is much larger than current
bounds suggest, or unless resonant enhancement factors ($Q_\psi$) are
much larger than expected.


%=============================================================================
% BIBLIOGRAPHY
%=============================================================================

\begin{thebibliography}{99}

\bibitem{Alcock2025DFD}
G.~T.~Alcock, ``Density Field Dynamics: A Complete Unified Theory,'' (2025).

\bibitem{Alcock2025EMcoupling}
G.~T.~Alcock, ``Universal Electromagnetic Actuation of Scalar Refractive
Fields: Theory, Experimental Evidence, and Applications,'' (2025).

\bibitem{Alcock2025DeepEM}
G.~T.~Alcock, ``Deep Analysis of Electromagnetic Coupling to Scalar
Refractive Fields,'' (2025).

\bibitem{Alcock2025Generator}
G.~T.~Alcock, ``Artificial Generation and Control of Scalar Density
Fields: Laboratory Apparatus Design,'' (2025).

\bibitem{Alcock2025StrongFields}
G.~T.~Alcock, ``Strong Fields and Gravitational Waves in Density Field
Dynamics,'' (2025).

\bibitem{BekensteinMilgrom1984}
J.~Bekenstein and M.~Milgrom, ``Does the missing mass problem signal
the breakdown of Newtonian gravity?'' Astrophys.\ J.\ \textbf{286},
7 (1984).

\bibitem{Einstein1911}
A.~Einstein, ``On the influence of gravitation on the propagation of
light,'' Ann.\ Phys.\ \textbf{35}, 898 (1911).

\bibitem{Einstein1912}
A.~Einstein, ``The speed of light and the statics of the gravitational
field,'' Ann.\ Phys.\ \textbf{38}, 355 (1912).

\bibitem{Gordon1923}
W.~Gordon, ``Zur Lichtfortpflanzung nach der Relativit\"atstheorie,''
Ann.\ Phys.\ \textbf{72}, 421 (1923).

\bibitem{Abbott2017GW}
B.~P.~Abbott \textit{et al.} (LIGO/Virgo/Fermi/INTEGRAL),
``Gravitational Waves and Gamma-Rays from a Binary Neutron Star Merger:
GW170817 and GRB~170817A,'' Astrophys.\ J.\ Lett.\ \textbf{848}, L13
(2017).

\bibitem{Bellini2014}
E.~Bellini and I.~Sawicki, ``Maximal freedom at minimum cost:
linear large-scale structure in general modifications of gravity,''
JCAP \textbf{1407}, 050 (2014).

\end{thebibliography}

\end{document}
